\section{Introduction}


The Non-Fungible Token (NFT) market has experienced explosive growth, with trading volumes exceeding \$60 billion in 2021-2022. However, this growth has led to significant fragmentation across chains, marketplaces, and liquidity pools. Users face challenges including:

\begin{itemize}
\item \textbf{Liquidity Fragmentation}: NFT liquidity is scattered across dozens of marketplaces on multiple chains
\item \textbf{Inefficient Price Discovery}: Users must manually check multiple platforms to find best prices
\item \textbf{Cross-Chain Complexity}: Moving NFTs between chains requires technical knowledge and multiple transactions
\item \textbf{Poor User Experience}: Managing wallets, tokens, and transactions across different platforms
\end{itemize}

Lux.market addresses these challenges through a unified marketplace architecture built on the Reservoir Protocol \cite{reservoir2023}, providing:

\begin{equation}
P_{optimal} = \min_{m \in M} \left( P_m + F_m + G_m \right)
\end{equation}

where $P_m$ is the price on marketplace $m$, $F_m$ is the platform fee, and $G_m$ is the gas cost, across all marketplaces $M$ in the aggregation network.

\subsection{Market Landscape}

The NFT marketplace ecosystem consists of three primary categories:

\begin{enumerate}
\item \textbf{Native Marketplaces}: Platform-specific marketplaces (OpenSea, LooksRare, X2Y2)
\item \textbf{Aggregators}: Multi-marketplace interfaces (Gem, Genie, Reservoir)
\item \textbf{Specialized Markets}: Vertical-specific platforms (Art Blocks, SuperRare)
\end{enumerate}

Lux.market operates as a hybrid aggregator-marketplace, combining aggregation capabilities with native marketplace features and deep integration into the Lux ecosystem.

